% sec:adm-lapse-shift — Lapse, Shift, and the Spatial Metric
\section{Lapse, Shift, and the Spatial Metric}
\label{sec:adm-lapse-shift}

The foliation of \cref{sec:adm-foliation} equips spacetime with spatial
slices $\Sigma_t$, a unit normal $n^\mu$, and an induced metric
$\spatialmetric$ --- but it does not prescribe how spatial coordinate
labels connect across slices.  A point $P$ at coordinates $x^i$ on
$\Sigma_t$ is linked to $\Sigma_{t+\d t}$ by the coordinate time
vector $(\partial/\partial t)^\mu$, which need not be normal to the
slice.  The two fields that encode its deviation from the normal ---
the lapse, already encountered in the foliation, and the shift --- are
the gauge variables of the 3+1 decomposition, and together with the
spatial metric they reconstruct the full spacetime line element.
\coderef{ADMSpacetime}{typeclass}

\paragraph{The time vector.}

In adapted coordinates $(t, x^i)$, the coordinate time vector
$(\partial/\partial t)^\mu$ connects each point on $\Sigma_t$ to the
point with the same spatial labels on $\Sigma_{t+\d t}$.  Decomposing
it into components normal and tangent to $\Sigma_t$ gives
%
\begin{equation}
  \left(\frac{\partial}{\partial t}\right)^\mu
  = \lapse\,n^\mu + \beta^\mu \,,
  \label{eq:adm-time-vector}
\end{equation}
%
where $\lapse$ is the lapse function from \cref{sec:adm-foliation} and
$\beta^\mu$ is a vector tangent to $\Sigma_t$ called the \emph{shift
vector}.  The coefficient of $n^\mu$ is consistent with the definition
of the lapse: contracting both sides with $n_\mu$ and using
$n_\mu\,n^\mu = -1$ together with $n_\mu\,\beta^\mu = 0$ gives
$n_\mu\,(\partial/\partial t)^\mu = -\lapse$.  This matches the
definition $n_\mu = -\lapse\,\cd{\mu}t$ from
\cref{eq:adm-normal-covector}, since
$(\partial/\partial t)^\mu\,\cd{\mu}t = 1$ in adapted coordinates.

\paragraph{The shift vector.}

The shift $\beta^\mu$ is tangent to $\Sigma_t$ by construction:
%
\begin{equation}
  n_\mu\,\beta^\mu = 0 \,.
  \label{eq:adm-shift-tangent}
\end{equation}
%
In adapted coordinates its time component vanishes, leaving only the
three spatial components $\shift$.  The physical meaning is geometric:
an Eulerian observer at coordinates $x^i$ on $\Sigma_t$ moves along
$n^\mu$ and arrives on $\Sigma_{t+\d t}$ at spatial coordinates
$x^i - \beta^i\,\d t$, not $x^i$.  The minus sign arises because the
shift appears with a plus sign in \cref{eq:adm-time-vector}: the
observer follows $n^\mu$ from $(t, x^i)$ and reaches the same
\emph{spatial point} on the next slice, but the coordinate grid has
shifted by $+\beta^i\,\d t$, so the observer's new coordinate labels
are $x^i - \beta^i\,\d t$.  Setting $\shift = 0$ makes the coordinate
lines perpendicular to the slices --- the Eulerian observer's worldline
is then a coordinate line.
\physnote{Nonzero shift allows the coordinate grid to co-move with
dynamical features.  Binary black hole simulations use a corotating
shift to keep the punctures near fixed grid locations
(\cref{sec:bssn-gauge}).}

\begin{figure}[htbp]
  \centering
  % adm-lapse-shift.tex — Lapse, shift, and time vector decomposition
% Inputable TikZ fragment for fig:adm-lapse-shift
% Requires: tikz with calc, arrows.meta (loaded by ehbook.cls)
%
% Schematic spacetime diagram showing two neighbouring spatial slices
% Sigma_t and Sigma_{t+dt}.  From point P on the lower slice, three
% vectors illustrate the ADM decomposition:
%   - alpha n^mu (normal, vertical)   — EHJetBlue
%   - (d/dt)^mu  (time vector, tilted) — EHHorizonCrimson
%   - beta^i dt  (shift, tangential)   — EHAccretionGold
% The shift connects where the normal arrives (Q) to where the time
% vector arrives (R) on the upper slice.
%
% Sheet coordinates reuse the foliation diagram's perspective layout:
%   Front-to-back offset: dx = -0.7, dy = +1.5
%   Lower: FL=(0.5,0) FR=(7.0,0.20) BR=(6.3,1.70) BL=(-0.2,1.50)
%   Upper: FL=(0.5,2.50) FR=(7.0,3.40) BR=(6.3,4.90) BL=(-0.2,4.00)

\begin{tikzpicture}[
    >={Stealth[length=5pt]},
    font=\footnotesize,
  ]

  % ── Lower sheet (Σ_t) ─────────────────────────────────────────────

  \fill[EHMarginGray, opacity=0.06]
    plot[smooth, tension=0.6] coordinates {
      (0.5, 0.0) (2.5, -0.08) (4.5, 0.06) (7.0, 0.20)}
    -- (6.3, 1.70) -- (-0.2, 1.50) -- cycle;

  \draw[EHMarginGray, semithick]
    plot[smooth, tension=0.6] coordinates {
      (0.5, 0.0) (2.5, -0.08) (4.5, 0.06) (7.0, 0.20)};

  \draw[EHMarginGray, thin, opacity=0.4]
    (-0.2, 1.50) -- (6.3, 1.70);
  \draw[EHMarginGray, thin, opacity=0.4]
    (0.5, 0.0) -- (-0.2, 1.50)  (7.0, 0.20) -- (6.3, 1.70);

  % Coordinate gridlines — transverse
  \draw[EHMarginGray, very thin, opacity=0.15]
    (0.27, 0.50) -- (6.77, 0.70)
    (0.03, 1.00) -- (6.53, 1.20);

  % Coordinate gridlines — longitudinal
  \draw[EHMarginGray, very thin, opacity=0.15]
    (2.13, 0.05) -- (1.43, 1.55)
    (3.75, 0.10) -- (3.05, 1.60)
    (5.38, 0.15) -- (4.68, 1.65);

  \node[font=\small, text=EHMarginGray, anchor=west]
    at (7.15, 0.20) {$\Sigma_t$};


  % ── Upper sheet (Σ_{t+dt}) ────────────────────────────────────────

  \fill[EHMarginGray, opacity=0.06]
    plot[smooth, tension=0.6] coordinates {
      (0.5, 2.50) (2.5, 2.44) (4.5, 2.96) (7.0, 3.40)}
    -- (6.3, 4.90) -- (-0.2, 4.00) -- cycle;

  \draw[EHMarginGray, semithick]
    plot[smooth, tension=0.6] coordinates {
      (0.5, 2.50) (2.5, 2.44) (4.5, 2.96) (7.0, 3.40)};

  \draw[EHMarginGray, thin, opacity=0.4]
    (-0.2, 4.00) -- (6.3, 4.90);
  \draw[EHMarginGray, thin, opacity=0.4]
    (0.5, 2.50) -- (-0.2, 4.00)  (7.0, 3.40) -- (6.3, 4.90);

  % Gridlines — transverse
  \draw[EHMarginGray, very thin, opacity=0.15]
    (0.27, 3.00) -- (6.77, 3.90)
    (0.03, 3.50) -- (6.53, 4.40);

  % Gridlines — longitudinal
  \draw[EHMarginGray, very thin, opacity=0.15]
    (2.13, 2.73) -- (1.43, 4.23)
    (3.75, 2.95) -- (3.05, 4.45)
    (5.38, 3.18) -- (4.68, 4.68);

  \node[font=\small, text=EHMarginGray, anchor=west]
    at (7.15, 3.40) {$\Sigma_{t+\mathrm{d}t}$};


  % ── Point P on lower sheet ────────────────────────────────────────
  \coordinate (P) at (3.0, 0.55);

  % ── Points on upper sheet ─────────────────────────────────────────
  % Q: where the Eulerian observer arrives (along n^μ, directly above P)
  \coordinate (Q) at (3.0, 3.20);
  % R: where the coordinate time vector arrives (same labels x^i)
  \coordinate (R) at (4.8, 3.40);


  % ── Unit normal α n^μ (EHJetBlue) ────────────────────────────────
  \draw[->, EHJetBlue, very thick, shorten >=2pt]
    (P) -- (Q);
  \node[font=\scriptsize, text=EHJetBlue, anchor=east, inner sep=2pt]
    at (2.95, 1.90) {$\alpha\,n^\mu$};


  % ── Time vector (∂/∂t)^μ (EHHorizonCrimson) ──────────────────────
  \draw[->, EHHorizonCrimson, very thick, shorten >=2pt]
    (P) -- (R);
  \node[font=\scriptsize, text=EHHorizonCrimson, anchor=west, inner sep=2pt]
    at (4.20, 1.85) {$(\partial/\partial t)^\mu$};


  % ── Shift displacement β^i dt (EHAccretionGold) ──────────────────
  \draw[->, EHAccretionGold, very thick, shorten >=2pt]
    (Q) -- (R);
  \node[font=\scriptsize, text=EHAccretionGold, anchor=south, inner sep=2pt]
    at (3.9, 3.38) {$\beta^i\,\mathrm{d}t$};


  % ── Right-angle marker at Q ──────────────────────────────────────
  %
  % Shows orthogonality between n^μ (vertical) and β^μ (tangential).
  \draw[thin] (3.0, 3.05) -- (3.15, 3.07) -- (3.15, 3.20);


  % ── Point markers (drawn last, on top) ───────────────────────────
  \fill (P) circle (1.5pt);
  \node[font=\scriptsize, anchor=north east, inner sep=2pt]
    at (P) {$P$};

  \fill (Q) circle (1.5pt);
  \fill (R) circle (1.5pt);

\end{tikzpicture}

  \caption[Decomposition of the coordinate time vector into lapse and
    shift]{%
    Decomposition of the coordinate time vector into lapse and shift.
    From a point $P$ on $\Sigma_t$, the unit normal $\alpha\,n^\mu$
    (blue) reaches the next slice at the same \emph{spatial point},
    while the coordinate time vector
    $(\partial/\partial t)^\mu$ (red) reaches the point with the same
    coordinate labels $x^i$.  The shift $\beta^i\,\mathrm{d}t$ (gold)
    is the tangential displacement between these two endpoints on
    $\Sigma_{t+\mathrm{d}t}$, encoding how the coordinate grid slides
    between slices.  The right angle at their junction reflects the
    orthogonality $n_\mu\,\beta^\mu = 0$.}
  \label{fig:adm-lapse-shift}
\end{figure}

\paragraph{The ADM line element.}

The three ADM fields --- lapse, shift, and spatial metric --- determine
the full spacetime geometry.  The metric components in adapted
coordinates $(t, x^i)$ follow from the decomposition
\cref{eq:adm-time-vector}.  Expanding
$g(\partial_t, \partial_t) = g(\lapse\,n + \beta,\,\lapse\,n + \beta)$
and using $n_\mu\,n^\mu = -1$, $n_\mu\,\beta^\mu = 0$ gives
$g_{00} = -\lapse^2 + \beta_k\,\beta^k$.  The spatial block
$g_{ij} = \gamma_{ij}$ follows directly: both $\partial_i$ and
$\partial_j$ are tangent to $\Sigma_t$, so the induced metric agrees
with the spacetime metric on these vectors
(\cref{eq:adm-induced-metric}).  The cross term
$g_{0i} = g(\lapse\,n + \beta,\,\partial_i)
= \lapse\,n_\mu\,(\partial_i)^\mu + g(\beta,\,\partial_i)
= \gamma_{ij}\,\beta^j = \beta_i$,
since $n_\mu\,(\partial_i)^\mu = 0$ and both $\beta$ and $\partial_i$
are tangent to $\Sigma_t$.  Collecting these:
%
\begin{equation}
  g_{00} = -\lapse^2 + \beta_k\,\beta^k \,,\qquad
  g_{0i} = \beta_i \,,\qquad
  g_{ij} = \gamma_{ij} \,,
  \label{eq:adm-metric-components}
\end{equation}
%
where $\beta_i = \gamma_{ij}\,\beta^j$ is the shift with its index
lowered by the spatial metric.  Written as a completed square, the
line element takes a form that separates temporal and spatial
contributions:
%
\begin{equation}
  \d s^2 = -\lapse^2\,\d t^2
    + \gamma_{ij}\bigl(\d x^i + \beta^i\,\d t\bigr)
                  \bigl(\d x^j + \beta^j\,\d t\bigr) \,.
  \label{eq:adm-line-element}
\end{equation}
%
The first term encodes the proper-time lapse between slices: an
Eulerian observer measures $\d\tau = \lapse\,\d t$, so
$-\lapse^2\,\d t^2 = -\d\tau^2$.  The second term is the squared
spatial distance after accounting for the coordinate drift introduced
by the shift.  Expanding the product recovers
\cref{eq:adm-metric-components}.

\paragraph{The inverse metric and volume element.}

Inverting the $4 \times 4$ system \cref{eq:adm-metric-components} gives
the contravariant components
%
\begin{equation}
  g^{00} = -\frac{1}{\lapse^2} \,,\qquad
  g^{0i} = \frac{\beta^i}{\lapse^2} \,,\qquad
  g^{ij} = \inversespatialmetric - \frac{\beta^i\,\beta^j}{\lapse^2} \,.
  \label{eq:adm-inverse-metric}
\end{equation}
%
A direct check confirms $g^{\mu\alpha}\,g_{\alpha\nu}
= \delta^\mu{}_\nu$: for $\mu = \nu = 0$,
$g^{00}\,g_{00} + g^{0k}\,g_{k0}
= (-1/\lapse^2)(-\lapse^2 + \beta_l\,\beta^l)
  + (\beta^k/\lapse^2)\,\beta_k = 1$.
The mixed case $(\mu,\nu) = (0,i)$ also vanishes:
$g^{00}\,g_{0i} + g^{0k}\,g_{ki}
= (-1/\lapse^2)\,\beta_i + (\beta^k/\lapse^2)\,\gamma_{ki}
= 0$.
The spatial block $g^{ij}$ differs from $\inversespatialmetric$ by the
correction $-\beta^i\beta^j/\lapse^2$, which vanishes for zero shift
and encodes the tilt of the coordinate lines away from the normal.

Since $\spatialmetric$ is a Riemannian metric on the spacelike
hypersurface $\Sigma_t$, it is non-degenerate, and the Schur complement
of the spatial block gives the determinant of the four-metric:
%
\begin{equation}
  \detmetric = \lapse\,\sqrt{\gamma} \,,
  \label{eq:adm-determinant}
\end{equation}
%
where $\gamma = \det(\gamma_{ij})$.  The Schur complement of
the spatial block is $g_{00} - \beta_i\,\gamma^{ij}\,\beta_j
= -\lapse^2 + \beta_k\,\beta^k - \beta_k\,\beta^k = -\lapse^2$, so
$\det(g_{\mu\nu}) = -\lapse^2\,\gamma$.
The volume element thus splits into a temporal piece controlled by the
lapse and a spatial piece controlled by the induced metric ---
the evolution code never needs to compute a full $4 \times 4$
determinant but instead tracks $\lapse$ and $\gamma$ separately.
\implnote{The factorisation \cref{eq:adm-determinant} means the
integration measure $\detmetric\,\d^4 x$ decomposes into
$\lapse\,\sqrt{\gamma}\,\d t\,\d^3 x$, separating the gauge
variable $\lapse$ from the dynamical volume $\sqrt{\gamma}$.}

\paragraph{Gauge freedom.}

The lapse $\lapse$ and the three components of the shift $\shift$
constitute four freely specifiable functions --- precisely the number
of coordinate degrees of freedom in four-dimensional general relativity.
The spatial metric $\spatialmetric$ and the extrinsic curvature
$\extrinsic$ (\cref{sec:adm-extrinsic-curvature}) carry the dynamical
content; the lapse and shift are pure gauge.  Different choices lead to
different coordinate representations of the same physical spacetime.
The next section introduces the extrinsic curvature, which together
with $\spatialmetric$ constitutes the initial data for the ADM
evolution.
\warnnote{The gauge freedom is a feature, not a defect.  Geodesic
slicing ($\lapse = 1$, $\shift = 0$) hits the physical singularity in
finite coordinate time, crashing the simulation.  The gauge conditions
of \cref{sec:bssn-gauge} are designed to avoid this.}

\paragraph{Codebase decomposition.}

The \codemodule{Spacetime.Class} module defines the
\texttt{ADMSpacetime} typeclass that any foliated spacetime must
implement.  Its interface mirrors the three-field decomposition:
%
\begin{lstlisting}[language=Haskell]
class Spacetime s => ADMSpacetime s where
  lapse            :: s -> Point s -> Double
  shift            :: s -> Point s -> Vec3 'Contravariant
  spatialMetric    :: s -> Point s -> Sym3 'Covariant
  invSpatialMetric :: s -> Point s -> Sym3 'Contravariant
\end{lstlisting}
%
The lapse is a scalar, the shift a contravariant three-vector, and the
spatial metric a symmetric $3 \times 3$ covariant tensor.  The phantom
type parameter \texttt{'Covariant}/\texttt{'Contravariant} enforces
index variance at compile time, making it impossible to pass a covariant
shift where a contravariant one is expected.  From these three fields the
full four-metric \cref{eq:adm-line-element} is assembled whenever the
geodesic integrator or evolution code needs it.
